%% =========================================================================
\section{Exercise 2: Parallel Operation of Two Synchronous Generators}
%% =========================================================================

The simulation model is extended to include a second generator (Generator~B)
connected in parallel with Generator~A via a synchronisation breaker at
$t=\SI{5}{s}$. Generator~A remains connected to the initial \MW{250} load and
operates at its setpoint of approximately \pu{1.01}. Both generators have
identical \MVA{555} ratings, AVRs, and governor characteristics ($R=5\%$).

%% ── SYNCHRONISATION CONDITIONS ────────────────────────────────────────────
Before paralleling two generators, four conditions must be satisfied to avoid
damaging transient currents at breaker closure: (1)~equal voltage magnitude,
(2)~equal frequency, (3)~matching phase sequence, and (4)~zero (or small)
phase-angle difference across the breaker. In practice, a synchroscope or
automatic synchroniser verifies these conditions before issuing the close
command. In this simulation the breaker closes at a fixed time ($t=\SI{5}{s}$)
regardless of phase-angle alignment, so the initial transients observed in the
results partly reflect imperfect synchronisation. The two tasks below explore
what happens when condition~(2) is deliberately violated by setting Gen~B's
governor reference above or below Gen~A's.

%% ── SIMULINK MODEL DIAGRAM ───────────────────────────────────────────────
\subsection{Simulink Model}

\figref{fig:ex2_diagram} shows the Exercise~2 Simulink model. Two identical
synchronous generators are connected through a synchronisation breaker. Each
generator has its own AVR and governor subsystem with independently adjustable
speed reference setpoints. The synchronisation breaker closes at $t=\SI{5}{s}$,
coupling the two machines onto a common bus. Measurement blocks on each
generator capture frequency, terminal voltage, active power, and reactive
power for comparison.

\begin{figure}[H]
  \centering
  \includegraphics[width=0.75\textwidth]{EX2_figure.png}
  \caption{Exercise 2 Simulink model: two identical synchronous generators
    connected in parallel through a synchronisation breaker. Independent AVR
    and governor subsystems allow setpoint variation for Tasks~3 and~4.}
  \label{fig:ex2_diagram}
\end{figure}

%% ── TASK 3 ────────────────────────────────────────────────────────────────
\subsection{Task 3 -- Generator B at Lower Frequency (0.98\,pu Setpoint)}

Generator~B's governor speed reference is set to \pu{0.98}, meaning it rotates
at a frequency approximately 2\% below Generator~A before synchronisation.

\subsubsection{Frequency (Task 3)}

\begin{figure}[H]
  \centering
  \includegraphics[width=0.65\textwidth]{task3_frequency_variation.png}
  \caption{Task 3: Frequency variation when Generator~B is paralleled at
    $f_\text{set} = \pu{0.98}$. Breaker closes at $t = \SI{5}{s}$.}
  \label{fig:task3_freq}
\end{figure}

Before the breaker closes at $t=\SI{5}{s}$, Gen~A operates at approximately
\pu{1.01} and Gen~B at \pu{0.98}. Upon breaker closure, synchronising torques
act between the two machines, causing transient frequency oscillations as the
rotors exchange synchronising power:

\begin{equation}
  P_\text{sync} = \frac{|V_1||V_2|}{X_s}\sin(\delta_1 - \delta_2)
  \label{eq:sync_power}
\end{equation}

The final system frequency is determined by the combined droop characteristic:

\begin{equation}
  \Delta P_L = -\left(\frac{1}{R_1} + \frac{1}{R_2}\right)\Delta f
  \label{eq:combined_droop}
\end{equation}

With $R_1 = R_2 = 0.05$ and equal machine ratings, the settled system
frequency is the arithmetic mean of the two setpoints:

\begin{equation}
  f_\text{sys,ss} \approx \frac{f_{nl,A} + f_{nl,B}}{2}
  = \frac{1.01 + 0.98}{2} = 0.995\,\text{pu}
  \label{eq:task3_freq}
\end{equation}

Accounting for the connected load and droop dynamics, the simulation confirms
a settled frequency of approximately \pu{1.008}.

\subsubsection{Active Power (Task 3)}

\begin{figure}[H]
  \centering
  \includegraphics[width=0.65\textwidth]{task3_active_power.png}
  \caption{Task 3: Active power sharing when Generator~B is paralleled at
    $f_\text{set} = \pu{0.98}$. Gen~B absorbs power (motoring condition).}
  \label{fig:task3_active}
\end{figure}

Active power sharing depends on where each generator's governor setpoint sits
relative to the system frequency. Applying \eqnref{eq:droop_power} to Gen~B:

\begin{equation}
  P_B = \frac{f_{nl,B} - f_\text{sys}}{S_P}
  = \frac{0.98 - 1.008}{S_P} < 0
  \label{eq:task3_PB}
\end{equation}

The result is negative, meaning Generator~B absorbs active power from the
network --- it is running as a \textbf{synchronous motor}. Generator~A has to
supply the full connected load plus whatever Gen~B is drawing. In practice,
paralleling a machine with a lower frequency setpoint immediately puts extra
load on the running generator and is not a safe way to synchronise.

\subsubsection{Voltage and Reactive Power (Task 3)}

\begin{figure}[H]
  \centering
  \includegraphics[width=0.65\textwidth]{task3_voltage.png}
  \caption{Task 3: Terminal voltage variation. A minor dip occurs at
    $t=\SI{5}{s}$ upon breaker closure; both AVRs restore $V_t$ to \pu{1.0}.}
  \label{fig:task3_voltage}
\end{figure}

\begin{figure}[H]
  \centering
  \includegraphics[width=0.65\textwidth]{task3_reactive_power.png}
  \caption{Task 3: Reactive power variation. Transient oscillations at
    $t=\SI{5}{s}$ settle once synchronism is established.}
  \label{fig:task3_reactive}
\end{figure}

A minor voltage dip occurs at $t=\SI{5}{s}$ as the breaker closes, caused by
the circulating current driven by the phase-angle difference between the two
machines. Both AVRs respond promptly and restore the terminal voltages to
\pu{1.0}. Reactive power undergoes transient oscillations corresponding to
the initial phase-angle mismatch, settling to steady values once synchronism
is established.

%% ── TASK 4 ────────────────────────────────────────────────────────────────
\subsection{Task 4 -- Generator B at Higher Frequency (1.02\,pu Setpoint)}

Generator~B's governor speed reference is raised to \pu{1.02}, exceeding
Generator~A's setpoint.

\subsubsection{Frequency (Task 4)}

\begin{figure}[H]
  \centering
  \includegraphics[width=0.65\textwidth]{task4_frequency_variation.png}
  \caption{Task 4: Frequency variation when Generator~B is paralleled at
    $f_\text{set} = \pu{1.02}$. Breaker closes at $t = \SI{5}{s}$.}
  \label{fig:task4_freq}
\end{figure}

Before connection, Gen~B runs faster than Gen~A. Upon breaker closure,
synchronising power surges from the faster Gen~B to the slower Gen~A
(following \eqnref{eq:sync_power}), causing Gen~B to decelerate and Gen~A to
accelerate until they lock at a common frequency. With $R_1 = R_2 = 0.05$
and equal machine ratings, the settled system frequency is the arithmetic
mean of the two setpoints:

\begin{equation}
  f_\text{sys,ss} \approx \frac{f_{nl,A} + f_{nl,B}}{2}
  = \frac{1.01 + 1.02}{2} = 1.015\,\text{pu}
  \label{eq:task4_freq}
\end{equation}

The final settled frequency is higher than in Task~3, reflecting the higher
combined setpoint energy of the system.

\subsubsection{Active Power (Task 4)}

\begin{figure}[H]
  \centering
  \includegraphics[width=0.65\textwidth]{task4_active_power.png}
  \caption{Task 4: Active power sharing when Generator~B is paralleled at
    $f_\text{set} = \pu{1.02}$. Gen~B supplies the majority of the load.}
  \label{fig:task4_active}
\end{figure}

With Gen~B's setpoint above the system frequency:

\begin{equation}
  P_B = \frac{f_{nl,B} - f_\text{sys}}{S_P}
  = \frac{1.02 - f_\text{sys}}{S_P} > 0
  \label{eq:task4_PB}
\end{equation}

Generator~B contributes positive active power, taking the majority of the load.
This is the desired operational outcome of synchronisation, and contrasts
starkly with Task~3.

\subsubsection{Voltage and Reactive Power (Task 4)}

\begin{figure}[H]
  \centering
  \includegraphics[width=0.65\textwidth]{task4_voltage.png}
  \caption{Task 4: Terminal voltage variation. A transient spike occurs at
    $t=\SI{5}{s}$; AVRs restore $V_t$ to \pu{1.0}.}
  \label{fig:task4_voltage}
\end{figure}

\begin{figure}[H]
  \centering
  \includegraphics[width=0.65\textwidth]{task4_reactive_power.png}
  \caption{Task 4: Reactive power variation after synchronisation at
    $f_\text{set} = \pu{1.02}$.}
  \label{fig:task4_reactive}
\end{figure}

A transient voltage spike occurs at $t=\SI{5}{s}$ (as opposed to a dip in
Task~3), because Gen~B was ahead in phase angle when connected. The
AVRs restore both terminal voltages to \pu{1.0} rapidly.

%% ── TASK 3 vs TASK 4 COMPARISON ──────────────────────────────────────────
\subsection{Task 3 vs Task 4 -- Direct Comparison}

The comparison of Tasks~3 and~4 highlights the dramatic effect of the governor
frequency setpoint on power sharing behaviour.

In Task~3 ($f_{\text{set},B} = \pu{0.98}$), Generator~B operates below the
system frequency and absorbs power --- a hazardous motoring condition. In
Task~4 ($f_{\text{set},B} = \pu{1.02}$), Generator~B operates above the
system frequency and assumes the majority of the system load --- the desired
operational outcome.
